% !TeX root = ../教材题库_全部条目_转写工作稿.tex
% Source: XJ-B1-C06-PR-P247-Y0922
% 第6章 / 6.4 用样本估计总体 / 教材 p.247 / PDF 第254页
% OCR未校对：公式与文字均可能有误；未完成数学TeX转写；禁止直接发给学生
% 原题及相邻原解截图（核对用）：../题目原图/XJ-B1-C06-PR-P247-Y0922-01.webp
% 本文件只在转写完成后作为题干使用；OCR 条目尚不是可靠题干。
\begin{bankitem}{p.247 练习 第3题}{第6章 / 6.4 用样本估计总体 / 教材 p.247 / PDF 第254页}{XJ-B1-C06-PR-P247-Y0922}{OCR未校对：公式与文字均可能有误；未完成数学TeX转写；禁止直接发给学生}
\begin{ocrdraft}
[教材p.247 / PDF第254页]
\par
3．为了监控一种零件的一条生产线的生产过程，检验员每天从该生产线上
\par
随机抽取10个零件，并测量其尺寸（单位：(m)，结果如下：
\par
97
\par
98
\par
102
\par
105
\par
107
\par
108
\par
109
\par
113
\par
114
\par
（D计算该零件抽样尺寸的极差．
\par
（2）计算该零件抽样尺寸的样本均值样本方差、2和样本标准差s.
\par
（3）将样本均值作为总体均值萨的估计值，样本标准差5作为总体标准差。
\par
的估计值，根据生产经验，在一天的抽检零件中，如果出现了尺寸在一3。，
\par
萨十3刁之外的零件，就认为这条生产线在这一天的生产过程中可能出现了异常情
\par
况，需对当天的生产过程进行检查．试利用估计值判断是否需对当天的生产过程
\par
进行检查．
\end{ocrdraft}
\end{bankitem}
